\documentclass[10pt,a4paper]{amsart}
\usepackage{amsmath,amssymb,amsthm,graphicx}
\usepackage{verbatim}
\newenvironment{reference}{\comment}{\endcomment}
%\newenvironment{reference}{}{}
\newenvironment{slogan}{\comment}{\endcomment}
\newenvironment{history}{\comment}{\endcomment}
% For commutative diagrams we use Xy-pic
\usepackage[all]{xy}

% We use 2cell for 2-commutative diagrams.
\xyoption{2cell}
% Theorem environments.
%
\theoremstyle{plain}
\newtheorem{theorem}[subsection]{Theorem}
\newtheorem{proposition}[subsection]{Proposition}
\newtheorem{lemma}[subsection]{Lemma}

\theoremstyle{definition}
\newtheorem{definition}[subsection]{Definition}
\newtheorem{example}[subsection]{Example}
\newtheorem{exercise}[subsection]{Exercise}
\newtheorem{situation}[subsection]{Situation}

\theoremstyle{remark}
\newtheorem{remark}[subsection]{Remark}
\newtheorem{remarks}[subsection]{Remarks}

\numberwithin{equation}{subsection}

% Macros
%
\def\lim{\mathop{\mathrm{lim}}\nolimits}
\def\colim{\mathop{\mathrm{colim}}\nolimits}
\def\Spec{\mathop{\mathrm{Spec}}}
\def\Hom{\mathop{\mathrm{Hom}}\nolimits}
\def\Ext{\mathop{\mathrm{Ext}}\nolimits}
\def\SheafHom{\mathop{\mathcal{H}\!\mathit{om}}\nolimits}
\def\SheafExt{\mathop{\mathcal{E}\!\mathit{xt}}\nolimits}
\def\Sch{\mathit{Sch}}
\def\Mor{\mathop{\mathrm{Mor}}\nolimits}
\def\Ob{\mathop{\mathrm{Ob}}\nolimits}
\def\Sh{\mathop{\mathit{Sh}}\nolimits}
\def\NL{\mathop{N\!L}\nolimits}
\def\CH{\mathop{\mathrm{CH}}\nolimits}
\def\proetale{{pro\text{-}\acute{e}tale}}
\def\etale{{\acute{e}tale}}
\def\QCoh{\mathit{QCoh}}
\def\Ker{\mathop{\mathrm{Ker}}}
\def\Im{\mathop{\mathrm{Im}}}
\def\Coker{\mathop{\mathrm{Coker}}}
\def\Coim{\mathop{\mathrm{Coim}}}

% Boxtimes
%
\DeclareMathSymbol{\boxtimes}{\mathbin}{AMSa}{"02}

%
% Macros for moduli stacks/spaces
%
\def\QCohstack{\mathcal{QC}\!\mathit{oh}}
\def\Cohstack{\mathcal{C}\!\mathit{oh}}
\def\Spacesstack{\mathcal{S}\!\mathit{paces}}
\def\Quotfunctor{\mathrm{Quot}}
\def\Hilbfunctor{\mathrm{Hilb}}
\def\Curvesstack{\mathcal{C}\!\mathit{urves}}
\def\Polarizedstack{\mathcal{P}\!\mathit{olarized}}
\def\Complexesstack{\mathcal{C}\!\mathit{omplexes}}
% \Pic is the operator that assigns to X its picard group, usage \Pic(X)
% \Picardstack_{X/B} denotes the Picard stack of X over B
% \Picardfunctor_{X/B} denotes the Picard functor of X over B
\def\Pic{\mathop{\mathrm{Pic}}\nolimits}
\def\Picardstack{\mathcal{P}\!\mathit{ic}}
\def\Picardfunctor{\mathrm{Pic}}
\def\Deformationcategory{\mathcal{D}\!\mathit{ef}}
\usepackage{cleveref}
\crefname{originallemma}{original lemma}{original lemmas}
\Crefname{originallemma}{Original lemma}{Original lemmas}
\makeatletter
\newlabel{lemma-section-direct-sum-quasi-compact}{{lemma-section-direct-sum-quasi-compact}{source}{}{sourcepointer-1}{}}
\newlabel{lemma-section-direct-sum-quasi-compact@cref}{{[originallemma][0][]lemma-section-direct-sum-quasi-compact}{[1][0][]source}{}{}{}}
\makeatother
\crefname{originallemma}{original lemma}{original lemmas}
\Crefname{originallemma}{Original lemma}{Original lemmas}
\makeatletter
\newlabel{lemma-limits-colimits}{{lemma-limits-colimits}{source}{}{sourcepointer-2}{}}
\newlabel{lemma-limits-colimits@cref}{{[originallemma][0][]lemma-limits-colimits}{[1][0][]source}{}{}{}}
\makeatother
\crefname{originallemma}{original lemma}{original lemmas}
\Crefname{originallemma}{Original lemma}{Original lemmas}
\makeatletter
\newlabel{lemma-construct-quasi-coherent-sheaves}{{lemma-construct-quasi-coherent-sheaves}{source}{}{sourcepointer-3}{}}
\newlabel{lemma-construct-quasi-coherent-sheaves@cref}{{[originallemma][0][]lemma-construct-quasi-coherent-sheaves}{[1][0][]source}{}{}{}}
\makeatother
\title{A quasi-coherent module is locally associated to a module on a neighbourhood of any point with quasi-compact neighbourhoods}
\author{The Stacks Project Authors}
\date{stacks-project:modules.tex:local-snapshot-20261012; GNU FDL 1.2 or later}
\begin{document}
\maketitle
\noindent\textit{Editorial note.} All delivered mathematical text is the author's own, copied from the pinned Stacks Project chapter source sources/stacks-modules/pinned-original-source.tex; no statement or proof step was rewritten, reordered or removed.\par\medskip
\noindent\textit{Editorial note.} Preamble provenance: the Stacks Project keeps one single shared preamble for every chapter; this card takes it directly from the staged sources/stacks-modules/preamble.tex as three verbatim spans (lines 8-12, 14-18 and 154-217), so no derived file was written and nothing was added or redefined. The wrapper class and the base packages are supplied by the assembler.\par\medskip
\noindent\textit{Editorial note.} Bibliography provenance: neither the counted statement nor its proof contains a citation (this lane rescans the whole span before writing the card), so the card carries no bibliography and no entry was reconstructed.\par\medskip
\noindent\textit{Editorial note.} Selection provenance: this card was opened under the file-based control contract (CONTROL.md), with the mechanical pre-card gate run on its statement and proof interval before the card was written.\par\medskip
\noindent\textit{Allowed standard prerequisites (editorial).}\begin{itemize}
\item The description of a quasi-coherent module locally as the cokernel of a map between direct sums of copies of the structure sheaf. Used at the start of the counted proof (source lines 1476-1486). Provides the presentation whose map has to be turned into a matrix over the ring of sections of a neighbourhood.
\item The description of the sections of a direct sum over a quasi-compact open, and the second construction of the module associated to a module. Both are invoked inside the counted proof (source lines 1497-1498, 1512 and 1526-1528) as pointers to lemmas of this chapter. Converts the abstract map into a matrix of sections over a smaller neighbourhood and identifies the resulting sheaf with the given one.
\end{itemize}\medskip
\ifdefined\hypertarget\hypertarget{sourcepointer-1}{}\fi
\noindent\textit{Source reference.} Original lemma lemma-section-direct-sum-quasi-compact: sources/stacks-modules/pinned-original-source.tex, source lines 1498--1498. Navigation pointer only: the counted text refers to the Stacks item lemma-section-direct-sum-quasi-compact, which belongs to another chapter of the project and is not part of this card. The pointer carries the target's own label as its number because this chapter's staged snapshot does not contain the other chapter's tag table; its content is not reproduced here.\par\medskip
\ifdefined\hypertarget\hypertarget{sourcepointer-2}{}\fi
\noindent\textit{Source reference.} Original lemma lemma-limits-colimits: sources/stacks-modules/pinned-original-source.tex, source lines 534--534. Navigation pointer only: the counted text refers to the Stacks item lemma-limits-colimits, which belongs to another chapter of the project and is not part of this card. The pointer carries the target's own label as its number because this chapter's staged snapshot does not contain the other chapter's tag table; its content is not reproduced here.\par\medskip
\ifdefined\hypertarget\hypertarget{sourcepointer-3}{}\fi
\noindent\textit{Source reference.} Original lemma lemma-construct-quasi-coherent-sheaves: sources/stacks-modules/pinned-original-source.tex, source lines 1425--1425. Navigation pointer only: the counted text refers to the Stacks item lemma-construct-quasi-coherent-sheaves, which belongs to another chapter of the project and is not part of this card. The pointer carries the target's own label as its number because this chapter's staged snapshot does not contain the other chapter's tag table; its content is not reproduced here.\par\medskip
\begin{lemma}
\label{lemma-quasi-coherent-module}
Let $(X, \mathcal{O}_X)$ be a ringed space.
Let $x \in X$ be a point.
Assume that $x$ has a fundamental system of quasi-compact neighbourhoods.
Consider any quasi-coherent $\mathcal{O}_X$-module $\mathcal{F}$.
Then there exists an open neighbourhood $U$ of $x$
such that $\mathcal{F}|_U$ is isomorphic to the
sheaf of modules $\mathcal{F}_M$ on $(U, \mathcal{O}_U)$
associated to some $\Gamma(U, \mathcal{O}_U)$-module $M$.
\end{lemma}

\begin{proof}
First we may replace $X$ by an open neighbourhood of $x$
and assume that $\mathcal{F}$ is isomorphic to the
cokernel of a map
$$
\Psi :
\bigoplus\nolimits_{j \in J}
\mathcal{O}_X
\longrightarrow
\bigoplus\nolimits_{i \in I}
\mathcal{O}_X.
$$
The problem is that this map may not be given by
a ``matrix'', because the module of global sections of a direct sum
is in general different from the direct sum of the modules
of global sections.

\medskip\noindent
Let $x \in E \subset X$ be a quasi-compact
neighbourhood of $x$ (note: $E$ may not be open).
Let $x \in U \subset E$ be an open neighbourhood of
$x$ contained in $E$.
Next, we proceed as in the proof of
Lemma \ref{lemma-section-direct-sum-quasi-compact}.
For each $j \in J$ denote
$s_j \in \Gamma(X, \bigoplus\nolimits_{i \in I} \mathcal{O}_X)$
the image of the section $1$ in the summand $\mathcal{O}_X$
corresponding to $j$. There exists a finite collection of opens
$U_{jk}$, $k \in K_j$ such that $E \subset \bigcup_{k \in K_j} U_{jk}$
and such that each restriction $s_j|_{U_{jk}}$
is a finite sum $\sum_{i \in I_{jk}} f_{jki}$
with $I_{jk} \subset I$, and $f_{jki}$ in the summand
$\mathcal{O}_X$ corresponding to $i \in I$. Set
$I_j = \bigcup_{k \in K_j} I_{jk}$. This is a finite set.
Since $U \subset E \subset \bigcup_{k \in K_j} U_{jk}$
the section $s_j|_U$ is a section of the finite direct sum
$\bigoplus_{i \in I_j} \mathcal{O}_X$.
By Lemma \ref{lemma-limits-colimits}
we see that actually $s_j|_U$ is a sum
$\sum_{i \in I_j} f_{ij}$ and
$f_{ij} \in \mathcal{O}_X(U) = \Gamma(U, \mathcal{O}_U)$.

\medskip\noindent
At this point we can define a module $M$ as the cokernel of the map
$$
\bigoplus\nolimits_{j \in J}
\Gamma(U, \mathcal{O}_U)
\longrightarrow
\bigoplus\nolimits_{i \in I}
\Gamma(U, \mathcal{O}_U)
$$
with matrix given by the $(f_{ij})$. By construction (2) of
Lemma \ref{lemma-construct-quasi-coherent-sheaves} we see that
$\mathcal{F}_M$ has the same presentation as $\mathcal{F}|_U$
and therefore $\mathcal{F}_M \cong \mathcal{F}|_U$.
\end{proof}

\end{document}
